\section{Introduction}
\label{sec:intro}
\begin{figure}[t]
\centering
\resizebox{\columnwidth}{!}{%
\begin{tikzpicture}[font=\scriptsize, line join=round]
  \def\R{2.05}
  \colorlet{seriesTyped}{seriesTypedc}
  \colorlet{seriesLLM}{seriesLLMc}
  \draw[mutedc!25, line width=0.3pt] (0,0.25*\R) -- (0.25*\R,0) -- (0,-0.25*\R) -- (-0.25*\R,0) -- cycle;
  \draw[mutedc!25, line width=0.3pt] (0,0.50*\R) -- (0.50*\R,0) -- (0,-0.50*\R) -- (-0.50*\R,0) -- cycle;
  \draw[mutedc!25, line width=0.3pt] (0,0.75*\R) -- (0.75*\R,0) -- (0,-0.75*\R) -- (-0.75*\R,0) -- cycle;
  \draw[mutedc!45, line width=0.3pt] (0,1.00*\R) -- (1.00*\R,0) -- (0,-1.00*\R) -- (-1.00*\R,0) -- cycle;
  \node[font=\tiny, text=mutedc, fill=white, inner sep=0.6pt] at (-45:0.50*\R*0.7071) {50};
  \node[font=\tiny, text=mutedc, fill=white, inner sep=0.6pt] at (-45:1.00*\R*0.7071) {100};
  \foreach \a in {90,0,-90,180}{ \draw[mutedc!40, line width=0.3pt] (0,0) -- (\a:\R); }
  \node[anchor=south, align=center, font=\footnotesize] at (90:{\R+0.12}) {Intelligence};
  \node[anchor=west, align=center, font=\footnotesize] at (0:{\R+0.12}) {Calibration};
  \node[anchor=north, align=center, font=\footnotesize] at (-90:{\R+0.12}) {Speed};
  \node[anchor=east, align=center, font=\footnotesize] at (-180:{\R+0.12}) {Cost\\(higher = cheaper)};
  \draw[seriesTyped, line width=1.1pt, fill=seriesTyped, fill opacity=0.10] (90:0.5306*\R) -- (0:0.7634*\R) -- (-90:0.8327*\R) -- (-180:0.5197*\R) -- cycle;
  \fill[seriesTyped] (90:0.5306*\R) circle (1.5pt);
  \fill[seriesTyped] (0:0.7634*\R) circle (1.5pt);
  \fill[seriesTyped] (-90:0.8327*\R) circle (1.5pt);
  \fill[seriesTyped] (-180:0.5197*\R) circle (1.5pt);
  \draw[seriesLLM, line width=0.9pt, densely dotted] (90:0.9578*\R) -- (0:0.9198*\R) -- (-90:0.7366*\R) -- (-180:0.3689*\R) -- cycle;
  \fill[seriesLLM] ([shift={(-1.3pt,-1.3pt)}]90:0.9578*\R) rectangle ++(2.6pt,2.6pt);
  \fill[seriesLLM] ([shift={(-1.3pt,-1.3pt)}]0:0.9198*\R) rectangle ++(2.6pt,2.6pt);
  \fill[seriesLLM] ([shift={(-1.3pt,-1.3pt)}]-90:0.7366*\R) rectangle ++(2.6pt,2.6pt);
  \fill[seriesLLM] ([shift={(-1.3pt,-1.3pt)}]-180:0.3689*\R) rectangle ++(2.6pt,2.6pt);
  \draw[seriesLLM!55, line width=0.9pt, dash pattern=on 3pt off 1.5pt] (90:0.5447*\R) -- (0:0.5929*\R) -- (-90:0.8179*\R) -- (-180:0.2736*\R) -- cycle;
  \draw[seriesLLM!55, fill=white, line width=0.7pt] ([shift={(-1.3pt,-1.3pt)}]90:0.5447*\R) rectangle ++(2.6pt,2.6pt);
  \draw[seriesLLM!55, fill=white, line width=0.7pt] ([shift={(-1.3pt,-1.3pt)}]0:0.5929*\R) rectangle ++(2.6pt,2.6pt);
  \draw[seriesLLM!55, fill=white, line width=0.7pt] ([shift={(-1.3pt,-1.3pt)}]-90:0.8179*\R) rectangle ++(2.6pt,2.6pt);
  \draw[seriesLLM!55, fill=white, line width=0.7pt] ([shift={(-1.3pt,-1.3pt)}]-180:0.2736*\R) rectangle ++(2.6pt,2.6pt);
  \draw[mutedc, line width=0.9pt, dashed] (90:0.3180*\R) -- (0:0.5706*\R) -- (-90:0.7990*\R) -- (-180:0.7534*\R) -- cycle;
  \fill[mutedc] ([shift={(-1.5pt,-1.3pt)}]90:0.3180*\R) -- ++(3pt,0) -- ++(-1.5pt,2.8pt) -- cycle;
  \fill[mutedc] ([shift={(-1.5pt,-1.3pt)}]0:0.5706*\R) -- ++(3pt,0) -- ++(-1.5pt,2.8pt) -- cycle;
  \fill[mutedc] ([shift={(-1.5pt,-1.3pt)}]-90:0.7990*\R) -- ++(3pt,0) -- ++(-1.5pt,2.8pt) -- cycle;
  \fill[mutedc] ([shift={(-1.5pt,-1.3pt)}]-180:0.7534*\R) -- ++(3pt,0) -- ++(-1.5pt,2.8pt) -- cycle;
  \tikzset{legsample/.style={inner sep=0pt, minimum width=0.5cm, minimum height=2.8pt}}
  \node[legsample, anchor=west] (L0) at (-3.3,{-\R-0.75}) {};
  \draw[seriesTyped, line width=1.1pt] (L0.west) -- (L0.east); \fill[seriesTyped] (L0.center) circle (1.5pt);
  \node[right=2pt of L0, inner sep=1pt] (T0) {Jev 1.13.0};
  \node[legsample, right=10pt of T0] (L1) {};
  \draw[seriesLLM, line width=0.9pt, densely dotted] (L1.west) -- (L1.east); \fill[seriesLLM] ([shift={(-1.3pt,-1.3pt)}]L1.center) rectangle ++(2.6pt,2.6pt);
  \node[right=2pt of L1, inner sep=1pt] (T1) {GPT-6 Luna, low effort};
  \node[legsample, anchor=west] (L2) at (-3.3,{-\R-1.12}) {};
  \draw[seriesLLM!55, line width=0.9pt, dash pattern=on 3pt off 1.5pt] (L2.west) -- (L2.east); \draw[seriesLLM!55, fill=white, line width=0.7pt] ([shift={(-1.3pt,-1.3pt)}]L2.center) rectangle ++(2.6pt,2.6pt);
  \node[right=2pt of L2, inner sep=1pt] (T2) {Gemini 3.1 Flash-Lite};
  \node[legsample, right=10pt of T2] (L3) {};
  \draw[mutedc, line width=0.9pt, dashed] (L3.west) -- (L3.east); \fill[mutedc] ([shift={(-1.5pt,-1.3pt)}]L3.center) -- ++(3pt,0) -- ++(-1.5pt,2.8pt) -- cycle;
  \node[right=2pt of L3, inner sep=1pt] (T3) {OpenDecision};
\end{tikzpicture}%
}
\caption{Scores of four systems on JevBench v1.4.2.2~\cite{jevbench2026v1422}, from three classes: a typed decision model (Jev 1.13.0), two prompted LLMs (GPT-6 Luna at low reasoning effort and Gemini 3.1 Flash-Lite) and a zero-shot ModernBERT classifier (OpenDecision). Each axis is the benchmark's 0--100 score on the same decisions. Intelligence is chance-corrected accuracy across the benchmark's tiers, calibration measures how well stated probabilities match outcomes, speed maps latency to a score (100 at 0.1\,s, 80 at 1\,s), and cost maps the price per 1{,}000 decisions to a score, so that higher means cheaper. OpenDecision's speed and cost are the benchmark's estimates rather than measurements. The results file reports no uncertainty estimates. On this snapshot, no system is best on every axis. Against Gemini 3.1 Flash-Lite, the typed decision model scores 1.4 points lower on intelligence and higher on the other three axes. The speed scale is logarithmic: Jev's median latency is 0.65\,s, against 1.44\,s for GPT-6 Luna at low effort. We chose one system per class to compare the classes, not to rank systems. Three open-weight typed models rank above Jev on the benchmark's composite score.}
\label{fig:radar}
\end{figure}


As LLM systems take on tasks that involve several steps, they need to decide how to proceed as well as generate responses. An agent may need to choose which tool to call next and when to stop~\cite{yang2026toward}. The surrounding system may also choose which model should handle a request~\cite{varangotreille2025doing} or whether a proposed action is risky~\cite{liu2026evaluating}. Many of these decisions involve a known set of answers. For the software taking the next step, a choice among those answers can be more useful than an explanation, especially if it comes with a probability that helps determine whether to act or seek further review.

One way to obtain such decisions is to ask the generative LLM already handling the user's request to make them as well. For example, the system could ask that model whether to handle a request itself or pass it to another model. However, using a generative LLM for these decisions has three drawbacks: 1) the caller pays for generating answer tokens and, with reasoning models, for thinking tokens~\cite{sui2025stop}; 2) the returned text must be parsed before software can use it; and 3) any reported confidence needs to be checked for the model, task and output format. Token probabilities can be well calibrated in some formats~\cite{kadavath2022language}.

When the same decision is made repeatedly, a conventional classifier offers another option. It is fast and often accurate, but it is trained for a fixed task and label set. Changing the task or labels often requires adapting the classifier or its output head. For systems whose questions change, the work needed to adapt the decision method therefore matters alongside the cost of each call.

Typed decision models address this need by allowing the caller to specify the question, its answer type and its possible answers at call time. The caller supplies the context to be evaluated, called the \emph{state}, and the model returns a probability for each answer rather than generating text. TypeSafe AI offers this interface in Jev, which it describes as a System One model, after Kahneman's fast, intuitive System~1~\cite{kahneman2011thinking,typesafe2026docs}. We use \emph{typed decision model} for this class, also called Jev-style or System One decision models. Several open implementations expose the same interface~\cite{convai2026laya,palmer2026kev,bespoke2026nimble,mapika2026decider,lee2026semif,nokia2026anyjev}, and OpenAI has announced a related Decisions API in limited preview~\cite{openai2026decisions}.

With these options available, a developer needs to understand both how they differ from existing methods and when they are useful. Fig.~\ref{fig:radar} illustrates the practical comparison: on the reported benchmark snapshot, none of the four systems is best on every axis~\cite{jevbench2026v1422}. This motivates a closer look at the earlier research behind typed decision models, the work required to adapt them to a deployment, and the results reported so far.

Earlier research provides a starting point for this comparison. Typed decision interfaces bring together established methods for classification, reading answer probabilities directly from a model, checking those probabilities against observed outcomes, and deferring uncertain decisions. Tang and Zheng review 28 arXiv papers posted between September 19 and 24, 2026, and relate the interface to these methods, constrained generation, reranking and cascades that pass decisions between models~\cite{tang2026typed}. Two concurrent working drafts also review the area: one follows the development from classical classifiers to Jev, including routing and cascades~\cite{awesomejev2026decisions}, and the other examines calibration, deciding when to accept or defer an answer, and open implementations~\cite{luo2026jev}. Surveys of routing~\cite{varangotreille2025doing,moslem2026dynamic} and LLM judges~\cite{gu2024survey} provide further background on these related tasks. Building on this work, our survey makes three contributions:
\begin{enumerate}
\item \textbf{Connections to earlier research.} We explain these connections from two perspectives: 1) routing motivates adapting a system's decisions to change without retraining the decision model, and 2) classification and answer-probability methods help explain how typed decisions are produced (Section~\ref{sec:lineages}).
\item \textbf{Adaptation cost.} We examine the work needed when a system adds a new model, asks a new kind of question or changes a deployment requirement. We distinguish resources that must be acquired from those that can be reused (Section~\ref{sec:adaptation}).
\item \textbf{Evidence and practical guidance.} We bring together independent evaluations to explain where these models are useful and what deployment risks remain. We then give practitioners a procedure for comparing alternatives on their own tasks (Sections~\ref{sec:evidence}--\ref{sec:guidance}).
\end{enumerate}

\noindent\textbf{Scope and sources.} To connect the recent Jev literature with this earlier work, we searched arXiv, GitHub and official documentation for public sources through September 30, 2026\rev{, and cite a few later studies where they bear directly on a finding}. We combined a name for Jev or its interface (Jev, System One model, typed decision model, decision model) with a term for one surveyed family (LLM routing, cascades, zero-shot classification, policy-conditioned guards, rubric judges, label-token readout, calibration), then followed citations from the works found.

Because most work specific to typed decision models appeared in the preceding two weeks as preprints or repositories, the evidence is still developing. At this moment, Jev has no technical report yet, so our description follows TypeSafe's documentation~\cite{typesafe2026docs}. We distinguish documented interface properties from empirical findings and identify results supported by a single study. We call a study \emph{independent} when its authors neither developed nor funded the system it measures, and attribute results to the researchers who measured them.
